% !TEX program = lualatex
\documentclass[11pt,a4paper]{article}
\usepackage{icat-common}
\usepackage{hyperref}

%-------------------------------------------------
%  独自マクロ定義
%-------------------------------------------------
\newcommand{\lhagma}{\mathcal{L}}
\newcommand{\rupa}{A} % 色（対象）

%-------------------------------------------------
%  定理環境の設定
%-------------------------------------------------
\theoremstyle{definition}
\newtheorem{definition}{定義}[section]
\newtheorem{axiom}{公理}[section]
\newtheorem{theorem}{定理}[section]
\newtheorem{scholium}{補足}[section]

\begin{document}

\title{空論（Sunyata-theory）：絶対不動点の証明論的真理保存と境界生成の形式化}
\author{千葉将臣}
\date{2026-04-20}
\maketitle

\begin{abstract}
本稿は、空数学（Sunyata-Mathematics）の基底理論として「空論」を定式化する。空論では、自明な圏における恒等射 $id: \dfix \to \dfix \to \dfix$ を、ゲンツェン流シーケント計算における真理保存則 $\dfix \vdash \dfix$ と対応させる。さらに、境界生成演算子を導入する補強理論として「界論」を与え、世俗的対象（色） $\rupa$ の生成と、その四段階のメタ観測を通じた真理保存則への帰着を証明論的に記述する。最後に、チョナン派の他空思想における残余（lha gma）を直観主義論理の諸意味論により解釈し、残余記号 $\lhagma$ と絶対不動点 $\dfix$ の同一視が、構文論的・意味論的にどの条件のもとで正当化されるかを明示する。
\end{abstract}

\section{空論：絶対不動点と真理保存の公理系}

空論は、差別化または分割が導入される以前の勝義諦を、形式体系上の基底状態として記述する試みである。

\begin{definition}[絶対不動点 $\dfix$]
絶対不動点 $\dfix$ とは、対象間の区別を導入しない基底状態を表す記号である。圏論的には対象が一つである自明な圏 $\mathbf{1}$ に対応し、証明論的には外部前提を追加せずに保持される基礎的な真理値を表す。
\end{definition}

\begin{axiom}[空の自己同一性と真理保存則（空空）]
空論における基底状態は、外部対象を経由しない自己同一性として表される。この自己同一性は、圏論的には恒等射として、証明論（Gentzen流）においては Identity rule に対応する真理保存則として、以下のように記述される。
\[
\text{圏論的表現：} \quad id_{\dfix} : \dfix \to \dfix \to \dfix
\]
\[
\text{証明論的表現：} \quad \frac{\dfix \quad \dfix}{\dfix} \quad \text{あるいは} \quad \dfix \vdash \dfix
\]
この状態では、対象水準の差異に基づく推論ステップは導入されない。したがって、カット消去（Cut-elimination）の極限として、$\dfix$ から $\dfix$ への保存関係のみが残る。
\end{axiom}

\section{界論による補強：対象生成と証明論的帰着}

本節では、空論における真理保存則に対して、知覚上または構文上の境界が導入された場合の変化を界論として定式化する。

\begin{definition}[境界生成 $\metanegation$ と色の対象化]
絶対不動点 $\dfix$ に対して、非対称なメタ観測としての境界生成演算子 $\metanegation$ が作用すると、基底状態は対象として区別可能な形式を取る。本稿では、その対象化された形式を世俗的対象（色） $\rupa$ と表す。
\[
\metanegation(\dfix) \mweq \rupa
\]
\end{definition}

\begin{theorem}[界論による空の正当化と真理保存への帰着]
対象化された $\rupa$ は、演算子 $\metanegation$ の再帰的適用により、再び基底状態へ帰着される。
\[
\metanegation^4 \rupa \mweq \dfix
\]
証明論の観点では、四句分別に対応する四段階のメタ観測（色 $\to$ 受 $\to$ 想 $\to$ 行）の完結は、以下の推論規則として表される。
\[
\frac{\metanegation^4 \rupa \mweq \dfix}{\dfix \vdash \dfix}
\]
これは、境界生成によって導入された対象水準の差異が、メタ観測の完結により真理保存則へ帰着されることを意味する。したがって、対象化された構成は独立した実体としてではなく、$\dfix \vdash \dfix$ へ還元可能な派生的構成として扱われる。
\end{theorem}

\section{他空思想の証明論的解釈：二重四句分別と残余の形式化}

本節では、本論において定義した絶対不動点 $\dfix$ を、単なる虚無ではなく、世俗的対象の排他プロセスの極限においてなお消去されない「残余（lha gma：他空）」として解釈する。この解釈を、チベット仏教チョナン派のドルポパ（Dolpopa Sherab Gyaltsen, 1292--1361）の議論、および直観主義論理（Intuitionistic Logic: IL）の諸体系を用いて証明論的に基礎付ける。

ドルポパは『山法・了義海』（Jonang Mountain Doctrine）\cite{dolpopa}において、伝統的な中観派の四句分別（自空）を単層的な否定に限定せず、その否定によって得られた「空」の状態に対しても再度四句分別を適用する二重四句分別を展開した。この手続きは、世俗的対象を排除した後に残る勝義の真理、すなわち他空としての残余を記述するものと解釈できる。本稿では、この残余を $\lhagma$ と表記し、直観主義論理の構文論および意味論を用いて形式化する。

ここで $\dfix$ と $\lhagma$ の関係を明確にしておく。$\dfix$ は空論内部で導入された基底状態の記号であり、$\lhagma$ は二重四句分別および直観主義論理の各意味論において得られる残余の記号である。両者は定義上同一ではないが、(i) 対象水準の述語をすべて排除した後にも消去されないこと、(ii) 追加の対象的内容を持たず真理保存則としてのみ機能すること、(iii) 境界生成後の推論をカット消去の極限へ帰着させること、という三条件を満たす場合に同一視される。したがって本稿では、次の対応を作業仮説として採用する。
\[
\lhagma \cong \dfix, \qquad \text{したがって} \qquad \lhagma \vdash \lhagma \cong \dfix \vdash \dfix .
\]
この対応により、$\lhagma$ は他空思想における残余の記号、$\dfix$ は同じ構造を空論の公理系内で表す不動点記号として位置付けられる。


\subsection{直観主義論理における矛盾元の導出可能性}

直観主義論理において $\bot$ は、通常、矛盾または構成不可能性を表す定数として導入される。この意味では、$\bot$ は体系の語彙に属する基礎記号であり、任意の文脈なしに純粋な直観主義論理から定理として導出されるものではない。もし無条件に $\vdash \bot$ が成立するならば、体系そのものが不整合であることを意味する。

しかし、$\bot$ は単に外部から与えられる記号に留まらない。ある命題 $P$ とその否定 $\neg P$、すなわち $P \to \bot$ が同一の局所文脈に現れる場合、直観主義論理の含意除去により、次のように体系内で $\bot$ が導出される。
\[
\frac{P \qquad P \to \bot}{\bot}
\]
したがって、$\bot$ は無前提に導出される定理ではなく、矛盾的な境界条件が与えられた局所文脈において、体系内部の推論規則によって到達される特異点として理解される。

本稿で $\lhagma := \bot$ と記す場合、それは直観主義論理全体の不整合を主張するものではない。むしろ、対象水準の構成が $P$ と $\neg P$ の衝突として局所化されたとき、その衝突が体系内推論によって $\bot$ へ収束することを、残余 $\lhagma$ の一つの下方向の表現として解釈するものである。この意味で、$\bot$ は天下りに導入される基礎記号であると同時に、矛盾的境界のもとでは体系内から到達される形式的特異点でもある。

この二重性は、本稿における残余 $\lhagma$ の扱いと構造的に対応している。すなわち、$\lhagma$ もまた、他空思想において外部的に指定される残余としての側面と、二重四句分別および直観主義論理上の排他過程を通じて内部的に到達される極限としての側面を持つ。本稿では、この対応を厳密な同一性としてではなく、次のような作業上の同型関係として理解する。
\[
\bot_{\mathrm{ext}} \cong \bot_{\mathrm{int}}
\qquad \leadsto \qquad
\lhagma_{\mathrm{ext}} \cong \lhagma_{\mathrm{int}} .
\]
ここで左辺は、矛盾元における外部定義と内部導出の一致を表し、右辺は、他空思想における残余の外部的指定と証明論的帰着の対応を表す。この対応は米田の補題そのものを適用するものではないが、外部から指定された対象と、その対象へ向かう文脈的振る舞いによる内部的特徴付けが対応するという意味で、米田型の類比として読むことができる。

この観点から見ると、$\lhagma_{\mathrm{ext}}$ は残余が何を指すか、すなわち指示の側面を表し、$\lhagma_{\mathrm{int}}$ はその残余がどのような排他過程・推論過程を通じて特徴付けられるか、すなわち意味の側面を表すと解釈できる。したがって、$\lhagma_{\mathrm{ext}} \cong \lhagma_{\mathrm{int}}$ は、意味と指示を単純に同一視する主張ではなく、外部的に指定された残余の指示が、体系内部の意味構成によって回収されることを表す補助的な解釈である。

\subsection{直観主義論理の各体系における表現例}

残余 $\lhagma$ は、直観主義論理の複数の代表的なフレームワークにおいて、以下のように定式化できる。以下の列挙は網羅的分類ではなく、上記の三条件を満たす限りで、空論における絶対不動点 $\dfix$ の意味論的解釈を与える表現例である。

\begin{enumerate}[label=\arabic*)]
    \item \textbf{直観主義論理（中論の四句否定）} \\
    世俗の任意の命題 $P$ に対し、あらゆる自性（対象水準の指示）を構成不可能な領域として排他する、純粋な論理構文による記述。この整理は、龍樹『中論』における四句分別批判を証明論的記法へ移すための形式化である\cite{nagarjuna}。
    \[
    \lhagma := \neg(P) \wedge \neg(\neg P) \wedge \neg(P \wedge \neg P) \wedge \neg(\neg P \wedge \neg\neg P)
    \]

    \item \textbf{直観主義論理（ドルポパの二重四句分別）} \\
    対象を消去する第一段階の四句否定 $\mathcal{C}$（自空）に対し、その消去作用それ自体を高階のメタ視点から再度評価し、勝義の残余を記述する入れ子構造。
    \[
    \mathcal{C} := \neg(X) \wedge \neg(\neg X) \wedge \neg(X \wedge \neg X) \wedge \neg(\neg X \wedge \neg\neg X)
    \]
    \[
    \lhagma := \neg(\mathcal{C}) \wedge \neg(\neg \mathcal{C}) \wedge \neg(\mathcal{C} \wedge \neg \mathcal{C}) \wedge \neg(\neg \mathcal{C} \wedge \neg\neg \mathcal{C})
    \]

    \item \textbf{直観主義論理（BHK意味論）} \\
    A. ハイティングの直観主義的形式化において構築不可能性の極限（矛盾元）として用いられる算術的等式 $0=1$ に対する、メタレベルの翻訳（下方向の極限）\cite{heyting}。ここでの $\bot$ は無条件に証明される定理ではなく、矛盾的境界条件のもとで体系内推論により到達される特異点として扱う。
    \[
    \bot \equiv (0 = 1)
    \]
    \[
    \lhagma := \bot
    \]

    \item \textbf{直観主義論理（クリプキ意味論）} \\
    境界状態（$\neg P \wedge \neg\neg P$）を強制（forcing）する可能世界 $w$ において、直観主義的な情報の永続性（monotonicity）に基づき、任意の未来世界 $w' \ge w$ において保持される自己言及的不動点の準形式的記述。
    \[
    \lhagma := \{ w \mid (\bot \leftarrow w \Vdash (\neg P \wedge \neg\neg P) \to \forall w' \ge w, w' \Vdash \lhagma) \}
    \]
\end{enumerate}

\begin{scholium}[クリプキ意味論における再帰的定義について]
上記の定式化において $\lhagma$ の定義右辺に $\lhagma$ 自体が現れる。
標準的なクリプキ意味論では、充足関係 $\Vdash$ は命題の構造に関する
帰納法によって定義されるため、定義の整礎性（well-foundedness）が要求される。
すなわち、定義はより単純な構造への参照として完結しなければならず、
自己言及的な循環定義は標準的枠組みでは許容されない。

この禁止の理由は定義漏れではなく、充足関係の判定が停止しないという
構造的問題に起因する。したがって、上記の式は標準的クリプキ意味論における
定義ではなく、残余の自己言及性を示すための準形式的記法として読むべきである。

ただし不動点意味論・ゲーム意味論等の拡張では、
この種の循環を制御された形で扱う枠組みが存在する。
$\lhagma$ の自己言及的性格はその拡張が必要な対象として位置付けられるが、
当該拡張の正当化は本稿の範囲外であり、現時点では未確定として保留する。
\end{scholium}

\begin{enumerate}[label=\arabic*), resume]
    \item \textbf{直観主義論理（位相的意味論）} \\
    対象 $P$ とその二重否定閉包 $\neg\neg P$ の間に置かれる境界領域を含む、トポロジー空間 $\mathcal{O}(X)$ における最大の開集合としての記述。
    \[
    \lhagma := \bigcup \left\{ U \in \mathcal{O}(X) \;\middle|\; \emptyset \subseteq (\neg P \cap \neg\neg P) \subseteq U \subseteq \lhagma \right\}
    \]

    \item \textbf{直観主義論理（領域理論）} \\
    情報量が最小である底部元 $\bot$ を初期値とし、誤解除去の作用素 $F$ を単調増加関数として反復適用することで、上方向への最小不動点として得られる極限状態の記述。
    \[
    \lhagma := \bigsqcup_{n = 0}^{\infty} F^n(\bot)
    \]
\end{enumerate}

\subsection{証明論的帰着と空論の正当化}

これら直観主義論理の諸体系における定式化は、本論の第1節で定義した空論の絶対不動点 $\dfix$、すなわち真理保存則 $\dfix \vdash \dfix$ を、残余 $\lhagma$ の意味論的解釈として説明できることを示している。

また、直観主義論理は、ハイティング算術における $0=1$ のような矛盾元を局所化することで、体系全体の崩壊を避けながら、対象的内容を持たない残余を扱う形式的余地を保持している\cite{heyting}。この局所化により、$\bot$ は外部から導入される基礎記号でありながら、矛盾的境界条件のもとでは体系内推論によって到達される点として解釈される。この二重性は、$\lhagma$ を外部的残余として指定しつつ、同時に証明論的帰着の極限として扱う本稿の構成を支える論理的模型となる。

したがって、界論的なメタ観測演算子 $\metanegation$ の4回適用による収束 $\metanegation^4 \rupa \mweq \dfix$ は、爆発原理（ex falso quodlibet）による体系全体の崩壊を回避しつつ、推論をカット消去の極限へ帰着させる形式的手続きとして位置付けられる。

\section{結論}

本稿では、空論を $\dfix \vdash \dfix$ という真理保存則として定式化し、界論をその真理保存則に境界生成が導入された場合の補強理論として位置付けた。界論において生成される世俗的対象 $\rupa$ は、四段階のメタ観測を通じて $\metanegation^4 \rupa \mweq \dfix$ へ帰着される。

さらに、他空思想における残余 $\lhagma$ を直観主義論理の諸体系において解釈し、$\lhagma \cong \dfix$ という対応を導入した。この対応により、$\dfix$ は空論の内部記号、$\lhagma$ は排他プロセス後の残余を表す外部的・意味論的記号として区別されつつ、真理保存則の水準では同一の役割を果たすことが示される。また、$\lhagma$ の外部的指定と内部的特徴付けの対応は、残余の指示が体系内の意味構成によって回収されるという補助的解釈を与える。以上により、空論、界論、および他空思想の残余概念は、証明論的帰着という共通の形式のもとで統一的に記述される。


\begin{thebibliography}{9}
\bibitem{dolpopa}
Dol po pa Shes rab rgyal mtshan.
\textit{Ri chos nges don rgya mtsho}（『山法・了義海』）.
Buddhist Digital Resource Center, Work MW21871.
\url{https://library.bdrc.io/show/bdr:MW21871}

\bibitem{nagarjuna}
N\=ag\=arjuna.
\textit{M\=ulamadhyamakak\=arik\=a}（『中論』）.
本稿では、四句分別および自性批判の論理的構造を参照する基礎文献として用いる。

\bibitem{heyting}
A. Heyting.
\textit{Intuitionism: An Introduction}.
North-Holland, 1956.
本稿では、直観主義論理および矛盾元を用いる形式化の参照文献として用いる。
\end{thebibliography}

\end{document}
